
\section*{Core Concept}
\begin{itemize}
    \item Understanding rate as a ratio (quantity per unit)
    \item Distance--speed--time problems
    \item Work rate problems
    \item Unit conversions in rates
\end{itemize}

\section*{Key Rules}

\subsection*{Basic Rate Formula}
\begin{align*}
\text{Rate} = \frac{\text{Quantity}}{\text{Time}}
\end{align*}

\subsection*{Distance--Speed--Time}
\begin{align*}
d = rt, \quad r = \frac{d}{t}, \quad t = \frac{d}{r}
\end{align*}

\subsection*{Work Rate (Combined Work)}
\begin{align*}
\frac{1}{T} = \frac{1}{a} + \frac{1}{b}
\end{align*}

\section*{Common SAT Traps}
\begin{itemize}
    \item Mixing units (hours vs minutes)
    \item Using the wrong formula form
    \item Adding times instead of rates in work problems
    \item Forgetting total distance in round trips
    \item Not identifying what quantity is constant
\end{itemize}

\section*{Worked Examples}

\subsection*{Example 1}
A car travels 120 miles in 2 hours.
\[
r = \frac{120}{2} = 60 \text{ mph}
\]

\subsection*{Example 2}
How far in 3 hours at 50 mph?
\[
d = 50 \times 3 = 150
\]

\subsection*{Example 3 (Work Rate)}
A completes a job in 4 hours, B in 6 hours.
\[
\frac{1}{T} = \frac{1}{4} + \frac{1}{6} = \frac{5}{12}
\Rightarrow T = \frac{12}{5}
\]

\subsection*{Example 4 (Round Trip)}
A car travels 60 mph going and 40 mph returning. Average speed:

Let distance each way be 120:
\[
t_1 = \frac{120}{60} = 2, \quad t_2 = \frac{120}{40} = 3
\]

\[
\text{Average speed} = \frac{240}{5} = 48
\]

\section*{Practice Questions}

\subsection*{Easy}
\begin{enumerate}
    \item A car travels 100 miles in 2 hours. Find speed.
    \item How far in 5 hours at 20 mph?
\end{enumerate}

\subsection*{Medium}
\begin{enumerate}
    \setcounter{enumi}{2}
    \item A job takes 8 hours alone. What is its rate?
    \item Two workers take 6 hours and 3 hours. Together?
\end{enumerate}

\subsection*{Hard}
\begin{enumerate}
    \setcounter{enumi}{4}
    \item A car goes 80 mph and returns at 40 mph. Find average speed.
    \item Pipe A fills a tank in 4 hours, Pipe B in 12 hours. Together?
\end{enumerate}

\section*{Answers}
\begin{enumerate}
    \item 50 mph
    \item 100 miles
    \item $\frac{1}{8}$ job per hour
    \item 2 hours
    \item $\frac{160}{3}$ mph (53.33 mph)
    \item 3 hours
\end{enumerate}

\section*{Teaching Tips}
\begin{itemize}
    \item Always track units carefully
    \item Emphasize "rate means per"
    \item Use tables for distance, rate, time
    \item In work problems, add rates not times
\end{itemize}

\section*{Speed Strategy}
\begin{itemize}
    \item Memorize $d = rt$
    \item Treat total work as 1 job
    \item Use simple numbers for distances
    \item Convert units before solving
\end{itemize}
